\documentclass{article}

% packages
\usepackage{amsmath, amsthm, thmtools, amsfonts, amssymb, luacode, catchfile, tikzducks, hyperref, ifthen}
\ifcsname c@kobocompile\endcsname
	\usepackage[a5paper, total={1072pt, 1448pt}, margin=10pt, includeheadfoot]{geometry} % set page margins
\else
	\usepackage[a4paper, margin=50pt, includeheadfoot]{geometry}
\fi
\usepackage[shortlabels]{enumitem}
\usepackage[skip=3pt, indent=0pt]{parskip}

% language
\usepackage[bidi=basic, layout=tabular, provide=*]{babel}
\ifcsname c@english\endcsname
	\babelprovide[main, import]{english}
\else
	\babelprovide[main, import]{hebrew}
	\babelprovide{rl}
\fi
%\babelfont{rm}{Libertinus Serif}
\babelfont{rm}[Renderer=Harfbuzz]{Libertinus Serif}
\babelfont{sf}{Libertinus Sans}
\babelfont{tt}{Libertinus Mono}

% style
\AddToHook{cmd/section/before}{\clearpage}	% Add line break before section
\linespread{1.3}
\setcounter{secnumdepth}{0}		% Remove default number tags from sections, this won't do well with theorems
\AtBeginDocument{\setlength{\belowdisplayskip}{3pt}}
\AtBeginDocument{\setlength{\abovedisplayskip}{3pt}}
\graphicspath{ {../images/} }

% operators
\DeclareMathOperator\cis{cis}
\DeclareMathOperator\Sp{Sp}
\DeclareMathOperator\tr{tr}
\DeclareMathOperator\im{Im}
\DeclareMathOperator\re{Re}
\DeclareMathOperator\diag{diag}
\DeclareMathOperator*\lowlim{\underline{lim}}
\DeclareMathOperator*\uplim{\overline{lim}}
\DeclareMathOperator\rng{rng}
\DeclareMathOperator\Sym{Sym}
\DeclareMathOperator\Arg{Arg}
\DeclareMathOperator\Log{Log}
\DeclareMathOperator\dom{dom}
\DeclareMathOperator\supp{Supp}
\DeclareMathOperator\var{Var}
\DeclareMathOperator\cov{Cov}

% commands
%\renewcommand\qedsymbol{\textbf{מש''ל}}
%\renewcommand\qedsymbol{\fbox{\emoji{lizard}}}
\newcommand{\Aa}[0]{\mathcal{A}}
\newcommand{\Bb}[0]{\mathcal{B}}
\newcommand{\CC}[0]{\mathbb{C}}
\newcommand{\Cc}[0]{\mathcal{C}}
\newcommand{\EE}[0]{\mathbb{E}}
\newcommand{\FF}[0]{\mathbb{F}}
\newcommand{\Ff}[0]{\mathcal{F}}
\newcommand{\Ii}[0]{\mathcal{I}}
\newcommand{\Gg}[0]{\mathcal{G}}
\newcommand{\Ll}[0]{\mathcal{L}}
\newcommand{\Mm}[0]{\mathcal{M}}
\newcommand{\NN}[0]{\mathbb{N}}
\newcommand{\Nn}[0]{\mathcal{N}}
\newcommand{\PP}[0]{\mathbb{P}}
\newcommand{\Pp}[0]{\mathcal{P}}
\newcommand{\QQ}[0]{\mathbb{Q}}
\newcommand{\RR}[0]{\mathbb{R}}
\newcommand{\Rr}[0]{\mathcal{R}}
\newcommand{\Ss}[0]{\mathcal{S}}
\newcommand{\TT}[0]{\mathbb{T}}
\newcommand{\Uu}[0]{\mathcal{U}}
\newcommand{\Vv}[0]{\mathcal{V}}
\newcommand{\Ww}[0]{\mathcal{W}}
\newcommand{\ZZ}[0]{\mathbb{Z}}
\newcommand{\acts}[0]{\circlearrowright}
\newcommand{\explain}[2] {
	\begin{flalign*}
		 && \text{#2} && \text{#1}
	\end{flalign*}
}
\newcommand{\maketitleprint}[0]{ \begin{center}
	%\begin{tikzpicture}[scale=3]
	%	\duck[graduate=gray!20!black, tassel=red!70!black]
	%\end{tikzpicture}	
	\includegraphics[width=6cm]{cover}
\end{center}
}

% theorem commands
\newtheoremstyle{c_remark}
	{}	% Space above
	{}	% Space below
	{}% Body font
	{}	% Indent amount
	{\bfseries}	% Theorem head font
	{}	% Punctuation after theorem head
	{.5em}	% Space after theorem head
	{\thmname{#1}\thmnumber{ #2}\thmnote{ \normalfont{\text{(#3)}}}}	% head content
\newtheoremstyle{c_definition}
	{3pt}	% Space above
	{3pt}	% Space below
	{}% Body font
	{}	% Indent amount
	{\bfseries}	% Theorem head font
	{}	% Punctuation after theorem head
	{.5em}	% Space after theorem head
	{\thmname{#1}\thmnumber{ #2}\thmnote{ \normalfont{\text{(#3)}}}}	% head content
\newtheoremstyle{c_plain}
	{3pt}	% Space above
	{3pt}	% Space below
	{\itshape}% Body font
	{}	% Indent amount
	{\bfseries}	% Theorem head font
	{}	% Punctuation after theorem head
	{.5em}	% Space after theorem head
	{\thmname{#1}\thmnumber{ #2}\thmnote{ \text{(#3)}}}	% head content

\ifcsname c@english\endcsname
	\theoremstyle{plain}
	\newtheorem{theorem}{Theorem}[section]
	\newtheorem{lemma}[theorem]{Lemma}
	\newtheorem{proposition}[theorem]{Proposition}
	\newtheorem*{proposition*}{Proposition}
	%\newtheorem{corollary}[theorem]{אין חלופה עברית}

	\theoremstyle{definition}
	\newtheorem{definition}[theorem]{Definition}
	\newtheorem*{definition*}{Definition}
	\newtheorem{example}{Example}[section]
	\newtheorem{exercise}{Exercise}[section]

	\theoremstyle{remark}
	\newtheorem*{remark}{Remark}
	\newtheorem*{solution}{Solution}
	\newtheorem{conclusion}[theorem]{Conclusion}
	\newtheorem{notation}[theorem]{Notation}
\else
	\theoremstyle{c_plain}
	\newtheorem{theorem}{משפט}[section]
	\newtheorem{lemma}[theorem]{למה}
	\newtheorem{proposition}[theorem]{טענה}
	\newtheorem*{proposition*}{טענה}
	%\newtheorem{corollary}[theorem]{אין חלופה עברית}

	\theoremstyle{c_definition}
	\newtheorem{definition}[theorem]{הגדרה}
	\newtheorem*{definition*}{הגדרה}
	\newtheorem{example}{דוגמה}[section]
	\newtheorem{exercise}{תרגיל}[section]

	\theoremstyle{c_remark}
	\newtheorem*{remark}{הערה}
	\newtheorem*{solution}{פתרון}
	\newtheorem{conclusion}[theorem]{מסקנה}
	\newtheorem{notation}[theorem]{סימון}
\fi

% Questions related commands
\newcounter{question}
\setcounter{question}{1}
\newcounter{sub_question}
\setcounter{sub_question}{1}

\ifcsname c@english\endcsname
	\newcommand{\question}[1][0]{
		\ifthenelse{#1 = 0}{}{\setcounter{question}{#1}}
		\section{Question \arabic{question}}
		\addtocounter{question}{1}
		\setcounter{sub_question}{1}
	}

	\newcommand{\subquestion}[1][0]{
		\ifthenelse{#1 = 0}{}{\setcounter{sub_question}{#1}}
		\subsection{Part \alph{sub_question}}
		\addtocounter{sub_question}{1}
	}
\else
	\newcommand{\question}[1][0]{
		\ifthenelse{#1 = 0}{}{\setcounter{question}{#1}}
		\section{שאלה \arabic{question}}
		\addtocounter{question}{1}
		\setcounter{sub_question}{1}
	}

	\newcommand{\subquestion}[1][0]{
		\ifthenelse{#1 = 0}{}{\setcounter{sub_question}{#1}}
		\subsection{סעיף \localecounter{letters.gershayim}{sub_question}}
		\addtocounter{sub_question}{1}
	}
\fi

% import lua and start of document
\directlua{common = require ('../common')}

\GetEnv{AUTHOR}

% headers
\author{\AUTHOR}
\date\today

\title{אנליזה פונקציונלית --- סיכום}
\setcounter{secnumdepth}{2}

\usepackage{fancyhdr}
\pagestyle{fancy}
\renewcommand{\headrulewidth}{0pt}

\begin{document}
\maketitle
\maketitleprint[teal]

\tableofcontents

\section{שיעור 1 --- 26.3.2025}
\subsection{רקע}
אנליזה פונקציונלית היא כמו אלגברה לינארית.
בקורס זה נחקור מרחבים וקטוריים והעתקות עליהם, אבל על מרחבים מורכבים יותר והעתקות מורכבות יותר.
נתחיל בשאלה,
\begin{exercise}
	יהי $(X, \rho)$ מרחב מטרי כלשהו, ונניח ש־$A \subseteq X$.
	נניח גם ש־${(a_n)}_{n = 1}^\infty \subseteq A$. \\
	מהם התנאים ההכרחיים על $A$ כך ש־$(a_n)$ תכלול תת־סדרת קושי?
\end{exercise}
נעבור לדוגמה וטענות מאינפי 1 לרענן את זכרוננו.
\begin{example}
	המרחב המטרי הכי אינטואיטיבי הוא $X = \RR$ ו־$\rho(x, y) = |x - y|$.
\end{example}
\begin{proposition}
	תהי $A \subseteq \RR$ כך ש־$A$ חסומה, ותהי ${(a_n)}_{n = 1}^\infty \subseteq A$,
	אז יש ל־$(a_n)$ תת־סדרת קושי.
\end{proposition}
\begin{proof}
	$A \subseteq [-R, R]$ עבור $R \in \RR$.
	נתחיל בהגדרה של $\Delta_0 = A$ ולכן יש אינסוף, ולכן יש בקטע $\Delta_0$ אינסוף נקודות של הסדרה, וכן $|\Delta_0| = 2R$.
	נבחן את הקטעים החוצים את $\Delta_0$, הם $[-R, 0], [0, R]$, נבחר את זה מביניהם שמכיל אינסוף נקודות של $(a_n)$ להיות $\Delta_1$, וכך נמשיך ונגדיר סדרה $(\Delta_n)$.
	נובע שהסדרה הנתונה היא סדרה יורדת, במובן ש־$\Delta_0 \supset \Delta_1 \supset \Delta_2 \supset \dots$, מתקיים גם $|\Delta_n| = \frac{|\Delta_0|}{2^n}$ לכל $n \in \NN$, ובכל $\Delta_n$ יש אינסוף נקודות של $(a_n)$.
	נבחר $a_{n_1} \in \Delta_1$ וכך באופן כללי גם $a_{n_k} \in \Delta_k$, לכן נובע $|a_{n_k} - a_{n_l}| \le \frac{1}{2^k}$ עבור $l \ge k$.
	לכן נובע שאכן ישנה תת־סדרת קושי בסדרה $(a_n)$.
\end{proof}
\begin{remark}
	טענה זו נכונה גם כאשר מסתכלים על מרחב $(\RR^n, \rho)$ עבור $\rho(x, y) = \sqrt{\sum_{i = 1}^{n} {(x_i - y_i)}^2}$.
\end{remark}
\begin{definition}[מרחב נורמי]
	יהי $V$ מרחב וקטורי מעל $\FF$ עבור $\FF \in \{ \RR, \CC \}$, ותהי פונקציה $\lVert \cdot \rVert : V \to \RR_{\ge0}$ המקיימת,
	\begin{enumerate}
		\item $x = 0_V \iff \lVert x \rVert = 0$
		\item $\forall \alpha \in \FF, \lVert \alpha x \rVert = |\alpha| \cdot \lVert x \rVert$
		\item $\forall x, y \in V, \lVert x + y \rVert \le \lVert x \rVert + \lVert y \rVert$
	\end{enumerate}
	אז $(V, \lVert \cdot \rVert)$ יקרא מרחב נורמי עם נורמה $\lVert \cdot \rVert$.
\end{definition}
\begin{definition}[מרחב l2]
	נגדיר את הקבוצה $l_2 = \{ x = (x_1, \dots) \mid \forall k \in \NN, x_k \in \RR, \sum_{i = 1}^{\infty} x_i^2 < \infty \}$.
	נגדיר גם,
	\[
		\lVert x \rVert = {\left(\sum_{i = 1}^{\infty} x_i^2\right)}^{\frac{1}{2}}
	\]
	אז המרחב הנורמי $l_2$ הוא הקבוצה והנורמה הללו.
\end{definition}
נבחין כי עלינו להוכיח שזהו אכן מרחב נורמי לפי ההגדרה.
\begin{theorem}[אי־שוויון קושי־שווארץ]
	מתקיים,
	\[
		\sum_{i = 1}^{n} |x_i| \cdot |y_i|
		\le {\left( \sum_{i = 1}^{n} x_i^2 \right)}^{\frac{1}{2}} \cdot {\left( \sum_{i = 1}^{n} y_i^2 \right)}^{\frac{1}{2}}
	\]
\end{theorem}
\begin{notation}
	נסמן $\langle x, y \rangle = \sum_{i = 1}^{n}  x_i y_i$.
\end{notation}
\begin{proof}
	עבור $t \in \FF$ סקלר כלשהו,
	\[
		0
		\le \langle x + ty, x + ty \rangle
		= \langle x, x \rangle + 2t \langle x, y \rangle + \langle y, y \rangle t^2
	\]
	עובדה ידועה היא $At^2 + Bt + C \ge 0 \implies B^2 - 4AC \le 0$ ולכן,
	\[
		{\left( \sum_{i = 1}^{n} x_i y_i \right)}^2
		\le {\left( \sum_{i = 1}^{n} x_i^2 \right)} \cdot {\left( \sum_{i = 1}^{n} y_i^2 \right)}
	\]
	ולכן,
	\[
		\left\lvert \sum_{i = 1}^{n} x_i y_i \right\rvert
		\le {\left( \sum_{i = 1}^{n} x_i^2 \right)}^{1/2} \cdot {\left( \sum_{i = 1}^{n} y_i^2 \right)}^{1/2}
	\]
	ואם נגדיר $x_i' = |x_i$ וכן $y_i' = |y_i|$ אז מאי־השוויון הנתון נובע,
	\[
		\sum_{i = 1}^{n} |x_i'| \cdot |y_i'|
		\le {\left( \sum_{i = 1}^{n} x_i^2 \right)}^{\frac{1}{2}} \cdot {\left( \sum_{i = 1}^{n} y_i^2 \right)}^{\frac{1}{2}}
	\]
\end{proof}
נעבור להוכחת ההגדרה של $l_2$, כלומר ההוכחה שהנורמה שהגדרנו היא אכן נורמה.
\begin{proof}
	\begin{align*}
		{\lVert x + y \rVert}^2
		& = \sum_{i = 1}^{\infty} {(x_i + y_i)}^2 \\
		& = \sum_{i = 1}^{\infty} x_i^2 + 2 \sum_{i = 1}^{\infty} x_i y_i + \sum_{i = 1}^{\infty} y_i^2 \\
		& \le {\lVert x\rVert}^2 + 2 \lVert x \rVert \cdot \lVert y \rVert + {\lVert y \rVert}^2 \\
		& = {(\lVert x \rVert + \lVert y \rVert)}^2 \\
		& \implies \lVert x + y \rVert \le \lVert x \rVert + \lVert y \rVert
	\end{align*}
\end{proof}
עתה משקיבלנו ש־$l_2$ הוא אכן מרחב נורמי, נוכל לדון בתכונותיו.
\begin{example}
	במרחב $(l_2, \lVert \cdot \rVert)$ נגדיר את שפת כדור היחידה במרחב,
	\[
		S = \{ x \in l_2 \mid \lVert x \rVert = 1 \}
	\]
	נבחין כי $S$ קבוצה חסומה ב־$l_2$.
	נבחר ${(l_n)}_{n = 1}^\infty$ המוגדרת על־ידי $l_n = (0, \dots, 1, \dots)$ כאשר $l_n^n = 1, l_n^m = 0$ לכל $m \ne n$.
	כמובן מתקיים $\lVert l_n \rVert = 1$ לכל $n \in \NN$, ולכן סדרת הנקודות חסומה ב־$S$.
\end{example}
\begin{proposition}
	${(l_n)}_{n = 1}^\infty \subseteq l_2$ איינה כוללת תת־סדרת קושי.
\end{proposition}
\begin{proof}
	נבחין כי $\lVert l_n - l_m \rVert = \sqrt{2}$ לכל $n \ne m$.
\end{proof}
\begin{notation}[כדור]
	עבור מרחב מטרי $(X, \rho)$, נסמן $B_r(x) = \{ x \in X \mid \rho(x, x_0) < r \}$.
\end{notation}
\begin{definition}[קבוצה חסומה לחלוטין]
	יהי מרחב מטרי $(X, \rho)$ מרחב מטרי ותהי $A \subseteq X$, אז נאמר ש־$A$ חסומה לחלוטין (Totally bounded) אם לכל $\epsilon > 0$ קיים מספר סופי של נקודות $\{ x_1, \dots, x_N \} \subseteq X$,
	כך שמתקיים $A \subseteq \bigcup_{i = 1}^N B_\epsilon(x_i)$.
\end{definition}
מיד נראה שימוש בהגדרה זו במשפט, ובכך ניתן הצדקה להגדרה הלכאורה משונה הזאת.
\begin{theorem}[שקילות לחסימות לחלוטין]\label{totally_bounded_set_equivalecy_theorem}
	יהי מרחב מטרי $(X, \rho)$ ותהי $A \subseteq X$, אז התנאים הבאים שקולים,
	\begin{enumerate}
		\item $A$ חסומה לחלוטין.
		\item בכל סדרה של $A$ ניתן לבחור תת־סדרת קושי.
	\end{enumerate}
\end{theorem}
משפט זה הוא משפט חשוב ומרכזי, ועל הקורא לשנן את הוכחתו. את ההוכחה אומנם נראה בהרצאות הבאות, אך נראה עתה שימושים למשפט זה.
נעבור למשפט פחות חשוב ומרכזי,
\begin{theorem}[שקילות חסימות במרחבים האוקלידיים]
	נניח ש־$X = \RR^m$, וכן ש־$\rho(x, y) = \sqrt{\sum_{i = 1}^{m} {(x_i - y_i)}^2}$, אז אם $A \subseteq \RR^m$ חסומה, אז היא חסומה לחלוטין.
\end{theorem}
\begin{proof}
	נחסום את $A$ על־ידי קובייה מספיק גדולה, נחלק את הקובייה לתת־קוביות מספיק קטנות (ההצדקה מגיעה מאינפי 3), ונוכל לחסום כל קובייה כזו בכדור.
	נסמן $\{ x_i \} \subseteq \RR^m$ את מרכזי הקוביות ונקבל $A \subseteq \bigcup_{j = 1}^N B_\epsilon(x_j)$ מהגדרת החלוקה של הקובייה החוסמת.
\end{proof}
\begin{proposition}
	ב־$(l_2, \lVert \cdot \rVert)$ נגדיר את הקבוצה,
	\[
		\Pi = \{ x = (x_1, \dots) \in l_2 \mid \forall i \in \NN, |x_i| \le \frac{1}{2^{i - 1}} \}
	\]
	אם $x \in \Pi$ אז $\sum_{n = 1}^{\infty} x_n^2 < \infty$, ובהתאם בהכרח $\Pi \subseteq l_2$. \\
	הקבוצה $\Pi$ חסומה לחלוטין.
\end{proposition}
\begin{proof}
	תהי $(x_1, \dots) \in \Pi$, ונגדיר $x_n^* = (x_1, \dots, x_n, \dots, 0, 0, \dots)$.
	נגדיר גם $\Pi^*_n = \{ x = (x_1, \dots, x_n, 0, \dots) \mid |x_n| \le \frac{1}{2^{n - 1}} \}$.
	הקבוצה $\Pi_n^*$ חסומה לחלוטין, זאת שכן הקבוצה שקולה לקבוצה ב־$\RR^n$, ונבחין כי היא חסומה, ולכן ההוכחה שראינו קודם עודנה תקפה ובהתאם $\Pi_n^*$ חסומה לחלוטין. \\
	נבחין כי
	\[
		{\lVert x - x_n^* \rVert}^2
		= \sum_{i = n + 1}^{\infty} x_i^2
		\le \sum_{i = n + 1}^{\infty} \frac{1}{2^{2i - 2}}
		= \sum_{i = n + 1}^{\infty} \frac{4}{4^i}
		= \frac{1}{4^n} \cdot \frac{1}{1 - \frac{1}{4}}
		= \frac{1}{3 \cdot 4^{n - 1}}
	\]
	ולכן $\lVert x - x_n^* \rVert \le \frac{1}{2^{n - 1}}$.
	יהי $\epsilon > 0$, אז $\Pi_n^*$ חסומה לחלוטין ולכן קיימים $y^1, \dots, y^n \in l_2$ כך שמתקיים,
	\[
		\Pi_n^* \subseteq \bigcup_{i = 1}^N B_\epsilon(y^i)
	\]
	נניח ש־$x_n^* \in B_\epsilon(y^i)$ ונוכל לבחור $n$ כך שמתקיים $\lVert x - x_n^* \rVert < \epsilon$, אז
	\[
		\lVert x - y^i \rVert
		\le \lVert x - x_n^* \rVert + \lVert x_n^* - y^i \rVert
		< 2 \epsilon
	\]
	נובע ש־$\Pi \subseteq \bigcup_{i = 1}^N B_{2\epsilon}(y^i)$.
\end{proof}
נבחין כי עתה ראינו שב־$l_2$ במרחב נורמי יש קבוצות חסומות, זהו אכן מרחב מעניין.

\section{שיעור 2 --- 2.4.2025}
\subsection{חסימות לחלוטין}
נראה את הוכחתם של שני משפטים שמומלץ לזכור.
המשפט הראשון הוא משפט\ \ref{totally_bounded_set_equivalecy_theorem}, בקורס זה נקרא לו משפט האוסדורף, זאת למרות שזהו רק משפט חלקי למשפט המוכר כמשפט בשם זה.
נעבור להוכחה.
\begin{proof}
	נניח ש־$(X, \rho)$ מרחב מטרי וש־$A \subseteq X$ חסומה לחלוטין, לכן ניתן לכסות את הקבוצה $A$ על־ידי מספר סופי של כדורים.
	נניח ש־${\{ x_n \}}_{n = 1}^\infty \subseteq A$ ונבחר $\epsilon = 1$ התחלתי.
	מכאן נסיק שקיים כדור $B_{\epsilon = 1}^1$ הכולל אינסוף נקודות בסדרה.
	נגדיר $V^1 = A \cap B_{\epsilon = 1}^1$ ונסיק $\operatorname{diam}(V^1) = \sup_{x, y \in V^1} \rho(x, y) \le 2$.
	אז $V^1$ כולל מספר אינסופי של נקודות של $\{ x_n \}$.
	אין ספק ש־$V^1$ חסומה לחלוטין.
	נפעל עכשיו באופן דומה על $B_{\epsilon = 1}^1$, הקבוצה $V^1$ חסומה לחלוטין ולכן ניתן לכסות אותה על־ידי מספר סופי של כדורים עבור $\epsilon = \frac{1}{2}$.
	נבחר כדור שמכיל אינסוף נקודות של הסדרה ב־$V^1$, נסמנו $B_{\epsilon = \frac{1}{2}}^2$, ונגדיר גם $V^2 = V^1 \cap B_{\epsilon = \frac{1}{2}}^2$.
	הפעם $\operatorname{diam}(V^2) \le 1$ ולכן $V^2$ חסומה לחלוטין וכוללת מספר אינסופי של נקודות של $\{ x_n \}$.
	נחזור על תהליך זה אינסוף פעמים. \\
	בכתוצאה מהתהליך נקבל $V^1 \supset V^2 \supset \dots \supset V^k \supset \dots$ וכן ש־$\operatorname{diam}(V^k) \le \frac{2}{k}$,
	ואף ש־$V^k$ כולל אינסוף נקודות של $\{ x_n \}$.
	נבחר $x_{n_1} \in V^1, x_{n_2} \in V^2, \dots$ ונקבל תת־סדרה ${\{ x_{n_k} \}}_{k = 1}^\infty \subseteq A$ כך ש־$\rho(x_{n_k}, x_{n_{k + l}}) \le \frac{2}{k} \to 0$, זאת שכן $x_{n_k}, x_{n_{k + l}} \in V^k$.
	קיבלנו אם כן שתת־הסדרה היא קושי.

	נעבור לכיוון השני, נניח שלכל סדרה יש תת־סדרת קושי ב־$A$.
	נניח בשלילה כי $A$ אינה חסומה לחלוטין, לכן קיים $\epsilon > 0$ עבורו אין כיסוי סופי של כדורים.
	מספיק להוכיח כי ישנה סדרה ${\{ x_n \}}_{n = 1}^\infty \subseteq A$ שאינה כוללת תת־סדרת קושי.
	נבחר $x_1 \in A$.
	לכן נוכל להסיק שקיימת $x_2 \in A$ כך ש־$\rho(x_1, x_2) \ge \epsilon$.
	נמשיך כך להשתמש באי־החסימות עבור $\epsilon$ כדי לבנות סדרה של אינסוף נקודות כאלה, כלומר $\rho(x_n, x_m) \ge \epsilon$ לכל $n, m \in \NN$ כך ש־$n \ne m$.
	לסדרה $\{ x_n \}$ אין תת־סדרת קושי, בסתירה להנחה.
\end{proof}

\subsection{מרחבים מטריים חשובים}
\begin{definition}[מרחב הפונקציות הרציפות]
	נגדיר את המרחב המטרי $(C[a, b], \lVert \cdot \rVert_\infty)$ עבור $C[a, b] = \{ f : [a, b] \to \RR \mid f \text{ is continuous} \}$ ו־$\lVert f \rVert = \max_{x \in [a, b]} |f(x)|$.
	זהו מרחב נורמי.
\end{definition}
\begin{definition}[חסימות במידה אחידה]
	נניח ש־$\Phi \subseteq C[a, b]$ ונניח שקיים $K > 0$ כך שמתקיים $|\varphi(x)| \le K$ לכל $x \in [a, b]$ ולכל $\varphi \in \Phi$, כאשר $K$ אינו תלוי ב־$x, \varphi$.
	במקרה זה נאמר ש־$\Phi$ חסומה במידה אחידה.
\end{definition}
\begin{example}
	נגדיר $\Phi = {\{ \sin(nx) \}}_{n = 1}^\infty$, וידוע כי $|\sin(nx)| \le 1$, אז $\Phi$ חסומה לחלוטין.
\end{example}
\begin{example}
	נגדיר $f_n(x) = \frac{x^2}{x^2 + {(1 - nx)}^2}$ עבור $n \in \NN$, אז,
	\[
		\forall x \in [0, 1], n \in \NN,\ |f_n(x)| < 1
	\]
	ולכן נאמר ש־$\{ f_n \}$ חסומה במידה אחידה.
\end{example}
\begin{definition}[רציפות במידה אחידה]
	באנגלית Eqicontinuous family of functions.
	נניח ש־$\Phi \subseteq C[a, b]$ עבור כל $\epsilon > 0$ קיים $\delta = \delta(\epsilon)$ (כלומר ערך $\delta$ תלוי רק ב־$\epsilon$), כך שמתקיים,
	\[
		\forall x_1, x_2 \in [a, b], \varphi \in \Phi\ 
		|x_1 - x_2| \le \delta(\epsilon)
		\implies |\varphi(x_1) - \varphi(x_2)| \le \epsilon
	\]
	במקרה זה $\Phi$ נקראת רציפה במידה אחידה.
\end{definition}
\begin{example}
	נחזור לדוגמה האחרונה שלנו, ונבדוק אם היא רציפה במידה אחידה,
	\[
		|f_n(\frac{1}{n}) - f_n(0)| = 1
	\]
	ולכן $\{ f_n \}$ אינה רציפה במידה אחידה.
\end{example}
\begin{proposition}
	נניח ש־${\{ f_n \}}_{n = 1}^\infty \subseteq C[a, b]$.
	נניח שקיים $K > 0$ כך ש־$|f_n(x)| \le K$ עבור כל $x \in [a, b], n \in \NN$.
	נניח גם ש־$|f_n'(x)| \le K$.
	אז הקבוצה $\{ f_n \}$ חסומה במידה אחידה וגם רציפה במידה אחידה.
\end{proposition}
\begin{proof}
	נבחין כי מתקיים, $|f_n(x_1) - f_n(x_2)| \le |f'(y)| \cdot |x_1 - x_2| \le K |x_1 - x_2|$. \\
	לכן ניתן לבחור $\delta(\epsilon) = \frac{\epsilon}{K}$, והוא לא תלוי בפונקציות או בערכי $n$.
\end{proof}

\section{שיעור 3 --- 9.4.2025}

\subsection{תכונות מרחבי פונקציות}
\begin{theorem}[משפט ארצלה]
	נניח ש־$\Phi \subseteq (C[a, b], \lVert \cdot \rVert_\infty)$, אז התנאים הבאים שקולים,
	\begin{enumerate}
		\item לכל סדרה ${\{ f_n \}}_{n = 1}^\infty \subseteq \Phi$ קיימת תת־סדרת קושי.
			כלומר קיימת $\{ f_{n_k} \}$ כך ש־$\lVert f_{n_k} - f_{n_{k + l}} \rVert_\infty \xrightarrow{k \to \infty} 0$ עבור כל $l \in \NN$.
		\item $\Phi$ חסומה במידה אחידה ורציפה במידה אחידה.
	\end{enumerate}
\end{theorem}
\begin{proof}
	בכיוון הראשון נניח שלכל סדרה יש תת־סדרת קושי.
	ממשפט\ \ref{totally_bounded_set_equivalecy_theorem} נסיק ישירות ש־$\Phi$ חסומה לחלוטין.
	נבחר $\epsilon > 0$ ולכן $\Phi \subseteq \bigcup_{i = 1}^N B_\epsilon(f_i)$.
	תהי $\varphi \in \Phi$, אז קיים $1 \le i \le N$ כך ש־$\varphi \in B_\epsilon(f_i)$,
	\[
		\lVert \varphi \rVert_\infty
		= \lVert \varphi - f_i + f_i \rVert_\infty
		\le \lVert \varphi - f_i \rVert_\infty + \lVert f_i \rVert_\infty
		\le \epsilon + \lVert f_i \rVert_\infty
	\]
	מסדרות קושי נוכל להסיק שקיימים,
	\[
		\forall x \in [a, b],\ |f_1(x)| \le K_1, \dots, |f_N(x)| \le K_N
	\]
	ולכן נגדיר $K = \max\{ K_1, \dots, K_N \}$, לכן מתקיים $\lVert \varphi \rVert_\infty \le \epsilon + K$, נובע ש־$\Phi$ חסומה במידה אחידה.

	נעבור להוכחת רציפות במידה שווה. \\
	הפונקציות $f_1, \dots, f_N$ רציפות בקטע $[a, b]$, לכן רציפות בו במידה שווה, ונובע שקיימים $\delta_1(\epsilon), \ldots, \delta_N(\epsilon) > 0$ כך שמתקיים,
	\[
		\forall x, y \in [a, b],\ 
		|x - y| \le \delta_i(\epsilon)
		\implies |f_i(x) - f_i(y)| \le \epsilon
	\]
	נגדיר $\delta = \min\{ \delta_i \mid i \in \NN \}$.
	קיים $i \in \{1, \ldots, N\}$ כך ש־$\varphi \in B_\epsilon(f_i)$, לכן,
	\[
		|\varphi(x) - \varphi(y)|
		\le \overbrace{|\varphi(x) - f_i(x)|}^{\le \lVert \varphi - f_i\rVert_\infty} + \overbrace{|f_i(x) - f_i(y)|}^{\le \epsilon} + \overbrace{|f_i(y) - \varphi(y)|}^{\le \lVert \varphi - f_i\rVert_\infty}
	\]
	נניח גם ש־$|x - y| \le \delta(\epsilon)$ ולכן,
	\[
		\forall \epsilon > 0,\ \exists \delta = \delta(\epsilon),\ |x - y| \le \delta(\epsilon)
		\implies |\varphi(x) - \varphi(y)| \le 3 \epsilon
	\]
	כלומר, מצאנו רציפות במידה שווה.

	נעבור עתה לכיוון השני, נניח ש־$\Phi$ חסומה ורציפה במידה שווה.
	יהי $\epsilon > 0$ ו־$\delta(\epsilon) > 0$ כך שמתקיים,
	\[
		|x - y| \le \delta(\epsilon)
		\implies \forall \varphi \in \Phi,\ |\varphi(x) - \varphi(y)| \le \epsilon
	\]
	נגדיר את הסדרות כך ש־$y_{i + 1} - y_i \le \epsilon$ וכן סדרה כך ש־$x_{i + 1} - x_i \le \delta(\epsilon)$, ונגדיר גם $x_0 = a, x_n = b$ וכן $y_m = K, y_0 = -K$.
	ברור כי אכן קיימות סדרות סופיות כאלה, ונבחן את הנקודות האלה כמשרות חלוקה על החלק המתאים במישור.
	המטרה שלנו היא לחלק את הגרף של $\varphi$ תוך שימוש בתיבות שהגדרנו.
	נגדיר את הפונקציה $\psi$ כך שהיא עוברת דרך נקודות בתיבות הללו כך שהיא מקרבת את גרף $\phi$ אך קטנה ממנה תמיד, כלומר נבחר את החיתוכים $\varphi(x_i)$ ואת הנקודות $y_i$ הגדולות ביותר שמתחת לנקודות אלה.
	עתה נבדוק את $\lVert \varphi - \psi \rVert_\infty$ עבור $x \in [a, b]$,
	\[
		|\varphi(x) - \psi(x)|
		\le |\varphi(x) - \varphi(x_i)| + |\varphi(x_i) - \psi(x_i)| + |\psi(x_i) - \psi(x)|
		\le \epsilon + \epsilon + |\psi(x_i) - \psi(x_{i + 1})
		\le 2 \epsilon + 3 \epsilon
	\]
	כלומר קיבלנו ש־$\lVert \varphi - \psi \rVert_\infty \le 5 \epsilon$, כלומר ניתן לחסום $\Psi \subseteq \bigcup_{\psi \in \Gamma} B_{5 \epsilon}(\psi)$ עבור $\Gamma$ קבוצת המצולעים שעוברים דרך הנקודות ברשת שהגדרנו,
	כלומר זוהי קבוצה סופית.
\end{proof}
\begin{definition}[מרחב מטרי שלם]
	מרחב מטרי $(X, \rho)$ יקרא שלם אם כל סדרת קושי מתכנסת לנקודה במרחב המטרי.
\end{definition}
\begin{theorem}[שלמות מרחב הפונקציות הרציפות]
	המרחב $(C[a, b], \lVert \cdot \rVert_\infty)$ הוא מרחב מטרי שלם.
\end{theorem}
\begin{proof}
	תהי סדרה ${\{ f_n \}}_{n = 1}^\infty \subseteq C[a, b]$ ונניח כי סדרה זו היא סדרת קושי.
	כלומר
	\[
		\forall \epsilon > 0 \exists N = N(\epsilon) \in \NN,\ \forall n, m \ge N(\epsilon) \lVert f_n - f_m\rVert_\infty \le \epsilon
	\]
	נובע שלכל $x \in [a, b]$ גם $|f_n(x) - f_m(x)| \le \epsilon$, זאת מהגדרת הנורמה על מקסימום.
	אם נבחר $x \in [a, b]$ אז ${\{ f_n(x) \}}_{n = 1}^\infty \subseteq \RR$, כלומר זוהי סדרת ממשיים ומשלמות הממשיים והעובדה כי זוהי סדרת קושי נסיק שקיים $\lim_{n \to \infty} f_n(x) = y_x$.
	לכל $x$ נגדיר $f(x) = y_x$, כלומר נבנה פונקציה שמתקבלת מהנקודות הגבוליות של סדרת הפונקציות.
	כאשר $m \to \infty$ מתקיים,
	\[
		\forall \epsilon > 0, \exists N = N(\epsilon),\ \forall x \in [a, b],\ |f_n(x) - f(x)| \le \epsilon
	\]
	ולכן,
	\[
		\forall \epsilon > 0,
		\exists N = N(\epsilon),
		\forall n \ge N(\epsilon), \forall x \in [a, b], \max |f_n(x) - f(x)| \le \epsilon
	\]
	אז נובע ש־$\lim_{n \to \infty} \lVert f - f_n \rVert_\infty = 0$ כפי שרצינו.
\end{proof}
ניזכר במשפט שאנו כבר יודעים
\begin{theorem}[משפט ויירשטראס להתכנסות במידה שווה]
	אז אם ${\{ f_n \}}_{n = 1}^\infty C[a, b]$, ו־$f_n \rightrightarrows f$ (כלומר הסדרה מתכנסת במידה שווה) אז $f$ רציפה.
\end{theorem}
ונראה משפט שקשור אליו וחשוב לא פחות.
\begin{theorem}[שלמות l2]
	המרחב המטרי $(l_2, \lVert \cdot \rVert)$, שנזכיר שמוגדר על־ידי,
	\[
		l_2 = \left\{ x \in \RR^\NN \middle| \sum_{i = 1}^\infty x_i^2 < \infty \right\},
		\qquad \lVert x \rVert = {\left(\sum_{i = 1}^\infty x_i^2\right)}^{\frac{1}{2}}
	\]
	הוא מרחב מטרי שלם.
\end{theorem}
\begin{proof}
	תהי סדרה ${\{ x^n \}}_{n = 1}^\infty \subseteq l_2$, ונניח שזוהי סדרת קושי.
	אז אנו יודעים כי,
	\[
		\forall \epsilon > 0, \exists N = N(\epsilon),
		\forall n, m \ge N(\epsilon), \lVert x^n - x^m \rVert^2 \le \epsilon
		\implies \sum_{i = 1}^\infty {(x_i^n - x_i^m)}^2 \le \epsilon^2
	\]
	נשתמש בתנאי ההכרחי להתכנסות בהחלט של טורים ממשיים ונקבל שעובר כל $i \in \NN$ מתקיים,
	\[
		\forall \epsilon > 0, \exists N(\epsilon),\ 
		{(x_i^n - x_i^m)}^2 \le \epsilon^2
	\]
	אם נקבע $i$ אז נקבל סדרה ${\{ x_i^n \}}_{n = 1}^\infty \subseteq \RR$ סדרת קושי, ונגדיר $x_i = \lim_{n \to \infty} x_i^n$.
	נגדיר את הסדרה ${\{ x_i \}}_{i = 1}^\infty \subseteq \RR$.
	נקבל שמתקיים ${(x_i^n - x_i)}^2 \le \epsilon^2$ לכל $n > N(\epsilon)$ ולכל $i \in \NN$.
	נבחר $M$ כלשהו, אז מתקיים,
	\[
		\sum_{i = 1}^M {(x_i^n - x_i^m)}^2 \le \epsilon^2
	\]
	ונובע,
	\[
		\lim_{M \to \infty} \sum_{i = 1}^M {(x_i^n - x_i^m)}^2
		= \sum_{i = 1}^M {(x_i^n - x_i)}^2
		\le \epsilon^2
	\]
	אז מתקיים,
	\[
		\sum_{i = 1}^\infty {(x_i^n - x_i)}^2 \le \epsilon^2
	\]
	בהתאם מתקיים, $\lim_{n \to \infty} \lVert x^n - x\rVert^2 = 0$, נבדוק,
	\[
		\sum_{i = 1}^\infty x_i^2
		= \sum_{i = 1}^\infty {(x_i - x_i^n + x_i^n)}^2
		= 2 \sum_{i = 1}^\infty {(x_i - x_i^n)}^2 + 2 \sum_{i = 1}^\infty {(x_i^n)}^2
		< \infty
	\]
	כלומר מצאנו סדרה גבולית והוכחנו שהיא במרחב שלנו.
\end{proof}
\begin{conclusion}
	נניח ש־${\{ f_n \}}_{n = 1}^\infty \subseteq C[a, b]$ סדרה חסומה במידה שווה ורציפה במידה שווה,
	אז קיימת תת־סדרה ${\{ f_{n_k} \}}_{k = 1}^\infty \subseteq \{ f_n \}$ שמתכנסת במידה שווה לפונקציה $f \in C[a, b]$.
\end{conclusion}
\begin{theorem}[אנלוגי למשפט ארצלה ל־l2]
	נניח ש־$\Phi \subseteq l_2$, אז התנאים הבאים שקולים,
	\begin{enumerate}
		\item $\Phi$ חסומה לחלוטין
		\item \begin{enumerate}
				\item קיים $K > 0$ כך ש־$\lVert \varphi \rVert \le K$ לכל $\varphi \in \Phi$, כלומר $\Phi$ חסומה
				\item $\lim_{M \to \infty} \left( \sup_{x \in \Phi} \sum_{i = M}^\infty x_i^2 \right) = 0$
		\end{enumerate}
	\end{enumerate}
\end{theorem}
ננסה להבין את התנאי שהרגע הגדרנו,
\begin{example}
	נגדיר את $S \subseteq l_2$ על־ידי $S = \{ x \mid \lVert x \rVert = 1 \}$.
	ב־$S$ נמצאות הסדרות $e_n = (0, \ldots, 0, 1, 0, \ldots)$ כאשר $e_n^n = 1$ בלבד.
	בהתאם מתקיים $\sup_{x \in S} \sum_{i = M}^\infty x_i^2 = 1$, לכן התנאי השני לא מתקיים ובהתאם לא יתכן ש־$S$ חסומה לחלוטין.
\end{example}
\begin{example}
	נגדיר את $H = \{ x \in l_2 \mid \forall n \in \NN, |x_n| \le \frac{1}{2^{n - 1}} \}$.
	הפעם נקבל,
	\[
		\sum_{i = M}^\infty x_i^2
		= \sum_{i = M}^\infty \frac{1}{4^{i - 1}}
		= \frac{4}{4^M} \cdot \frac{1}{1 - \frac{1}{4}}
		\xrightarrow{M \to \infty} 0
	\]
	ולכן התנאי השני עבור חסימות לחלוטין מתקיים.
\end{example}

\section{ שיעור 4 --- 23.4.2025}
\subsection{תכונות מרחבי סדרות}
נסיים את הפרק הזה בתכונות חשובות במרחב $(l_2, \lVert \cdot \rVert)$ עליו דנו בשיעורים הקודמים.
\begin{theorem}[משפט ארצלה ל־l2]
	נניח ש־$K \subseteq l_2$, אז התנאים הבאים שקולים,
	\begin{enumerate}
		\item $K$ חסומה לחלוטין
		\item
			\begin{enumerate}
				\item הקבוצה $K$ חסומה במרחב המטרי $(l_2, \lVert \cdot \rVert)$, ו
				\item $\lim_{M \to \infty} \sup_{x \in K} \sum_{j = M}^\infty x_j^2 = 0$
			\end{enumerate}
	\end{enumerate}
\end{theorem}
לפני שניגש להוכחת המשפט, נוכיח טענה כללית שתעזור לנו.
\begin{proposition}
	נניח ש־$(X, \rho)$ מרחב ממטרי כלשהו ונניח ש־$Q \subseteq X$ חסומה לחלוטין.
	אז $Q$ היא חסומה ב־$(X, \rho)$.
\end{proposition}
\begin{proof}
	תהי נקודה $x_0 \in X$, $Q$ חסומה לחלוטין ולכן $Q \subseteq \bigcup_{i = 1}^N B_{\epsilon}(x_i)$ עבור $N \in \NN$ ו־$x_1, \ldots, x_N \in X$. \\
	נגדיר $R = \max\{ \rho(x_0, x_1), \ldots, \rho(x_0, x_N) \}$.
	אם $q \in Q$ אז $q \in B_{\epsilon}(x_i)$ עבור איזשהו $i$, נובע שגם,
	\[
		\rho(q, x_0)
		\le \rho(q, x_i) + \rho(x_i, x_0)
		\le \epsilon + R
	\] 
	לכל $q \in Q$, נסיק שמתקיים $\rho(q, x_0) \le R + \epsilon$.
\end{proof}
נעבור להוכחת המשפט.
\begin{proof}[הוכחת המשפט]
	$1 \implies 2$, טענה $(a)$ נובעת מיד מהטענה שהוכחנו זה עתה.
	נעבור להוכחת $(b)$.
	יהי $\epsilon > 0$, $K$ קבוצה חסומה לחלוטין, ב־$l_2$ ולכן,
	\[
		K
		\subseteq \bigcup_{i = 1}^N B_{\epsilon}(x^n)
	\] 
	נבחין כי,
	\[
		\sum_{j = 1}^\infty {(x_j^1)}^2 < \infty
		, \ldots, 
		\sum_{j = 1}^\infty {(x_j^N)}^2 < \infty
	\] 
	אז קיימים $M_1, \ldots, M_N$ התלויים ב־$\epsilon$ בלבד כך שמתקיים,
	\[
		\sum_{i = M_1}^\infty {(x_i^1)}^2 \le \epsilon
		, \ldots, 
		\sum_{i = M_N}^\infty {(x_i^N)}^2 \le \epsilon
	\] 
	עבור $x = (x_1, \ldots) \in K$ מתקיים $x \in B_{\epsilon}(x^n)$ וכן $\lVert x - x^n \rVert^2 \le \epsilon^2 \le 2 \epsilon$ ולכן,
	\[
		\sum_{i = M}^\infty x_i^2
		= \sum_{i = M}^\infty {(x_i - x_i^n + x_i^n)}^2
		\le 2 \sum_{i = M}^\infty {(x_i - x_i^n)}^2 + 2 \sum_{i = M}^\infty {(x_i^n)}^2
	\]
	אז,
	\[
		\forall \epsilon > 0,
		\exists M,
		\forall x \in K,
		\ \sum_{i = M}^\infty x_i^2
		\le 2 \epsilon^2 + 2 \epsilon
	\] 
	ולכן למעשה מצאנו שמתקיים,
	\[
		\lim_{M \to \infty} \sup_{x \in K} \sum_{i = M}^\infty x_i^2 = 0
	\] 

	$2 \implies 1$,
	נניח ש־$K$ חסומה וכן שקיים הגבול (b).
	יהי $\epsilon > 0$ ונבחר $M$ כך שמתקיים,
	\[
		\sup_{x \in K} \sum_{i = M}^\infty x_i^2
		\le \epsilon^2
	\] 
	ולכן בפרט לכל $x \in K$, מתקיים $\sum_{i = M}^\infty x_i^2 \le \epsilon^2$
	נגדיר $\pi_M : K \to \pi_M(K) \subseteq {(\RR^M)}^\circ$ על־ידי $\pi_M(x) = (x_1, \ldots, x_M, 0, \ldots)$.
	אז $\pi_M(K)$ חסומה ב־$\RR^M$ ולכן $\pi_M(K)$ חסומה לחלוטין ב־$\RR^M$.
	נעיר שבמקרה זה $\RR^M = \{ (x_1, \ldots, x_M, 0, \ldots) \}$. \\
	נובע שקיימים $y^1, \ldots, y^N \in {(\RR^M)}^\circ$ כך שמתקיים,
	\[
		\pi_M(K)
		\subseteq \bigcup_{n = 1}^N B_{\epsilon}(y^n)
	\]
	אז אם $x \in K$, מתקיים $\pi_M(x) \in B_{\epsilon}(y^n)$, נסיק,
	\[
		\lVert x - y^n \rVert^2
		= \sum_{i = 1}^M {(x - y_i^n)}^2 + \sum_{i = M + 1}^\infty x_i^2 
		\le \lVert \pi_M(x) - y^n \rVert^2 + \epsilon^2
		\le 2 \epsilon^2
	\]
	בהתאם נובע ש־$K \subseteq \bigcup_{n = 1}^N B_{\sqrt{2} \epsilon}(y^n)$.
\end{proof}

\subsection{קירובים}
בעולם של אנליזה פונקציונלית עלינו למצוא דרך לקרב פונקציות מורכבות על־ידי פונקציות פשוטות יותר, זאת כדי שנוכל לעבוד במרחבים ההרבה יותר מורכבים שבהם וקטור הוא פונקציה.
עוד משהו שחשוב שנוכל לעשות הוא לקרב במידה שווה את הפונקציות, זאת שכן קירוב נקודתי לא מספר לנו מספיק על הפונקציות.
\begin{theorem}[משפט הקירוב של ויירשטראס]
	לכל $f \in C[0, 1]$ קיימת סדרת פולינומים ${(P_n)}_{n = 1}^\infty$ כך שמתקיים $P_n \rightrightarrows f$.
\end{theorem}
\begin{proof}
	נתחיל ברידוד של הבעיה, נניח ש־$g(x) = f(x) - f(0) - x(f(1) - f(0))$, אז נובע ש־$f(x) = g(x) + f(0) + x(f(1) - f(0))$, אך החלק המוסף הוא פולינום, ולכן נוכל לבחון את הקירוב ל־$g$ בלבד.
	נקבל שנוכל להניח ללא הגבלת הכלליות ש־$f(0) = f(1) = 0$. \\
	נגדיר פונקציה חדשה,
	\[
		F(x)
		= \begin{cases}
			f(x) & x \in [0, 1] \\
			0 & x \in \RR \setminus [0, 1]
		\end{cases}
	\] 
	פונקציה זו מוגדרת על הממשיים והיא רציפה במידה שווה ב־$\RR$ בשל ההנחה שעשינו. \\
	לכל $\epsilon > 0$ קיים $\delta > 0$ כך שאם $|x - y| \le 2 \delta$ אז $|F(x) - F(y)| \le \epsilon$ לכל $x, y \in \RR$.
	בשלב הבא נגדיר את סדרת הפולינומים שלנו,
	\[
		P_n(x)
		= \int_{-1}^{1} F(x + u) Q_n(u)\ du
	\]
	כאשר $Q_n(u) = C_n {(1 - u^2)}^n$ כאשר $C_n$ קבוע מנרמל כך שיתקיים $\int_{-1}^{1} Q_n(u)\ du = 1$. \\
	נבחין כי $F(x + u) \ne 0$ כאשר $x + u \in [0, 1]$, או בהתאם כאשר $u \in [-x, 1 - x]$.
	נשתמש בהגדרת התומך ונקבל שמתקיים,
	\[
		P_n(x)
		= \int_{-x}^{1 - x} F(x + u) Q_n(u)\ du
		= \int_0^1 F(t) Q_n(t - x)\ dt
	\]
	אבל $Q_n$ פולינום ונסיק שגם $P_n$ פולינום (מדוע?).
	נבחין כי,
	\begin{align*}
		& |P_n(x) - F(x)| \\
		& = \left\lvert \int_{-1}^{1} F(x + u) Q_n(u)\ du - \int_{-1}^{1} F(x) Q_n(u)\ du \right\rvert \\
		& \le \int_{-1}^{1} |F(x + u) - F(x)| Q_n(u)\ du \\
		& \le \int_{-1}^{-\delta} |F(x + u) - F(x)| Q_n(u)\ du + \int_{-\delta}^{\delta} |F(x + u) - F(x)| Q_n(u)\ du + \int_{\delta}^1 |F(x + u) - F(x)| Q_n(u)\ du
	\end{align*}
	נסמן את שלושת האינטגרלים הללו $I_1, I_2, I_3$ בהתאמה, ונבדוק,
	\[
		I_2
		\le \epsilon \int_{-\delta}^{\delta} Q_n\ du
		\le \epsilon \int_{-1}^{1} Q_n(u)\ du
		\le \epsilon
	\]
	עבור $I_3$, אנו יודעים ש־$F$ חסומה ונסמן $|F(x)| \le M$ עבור $M > 0$ כלשהו, אז,
	\[
		I_3
		\le 2M \int_{\delta}^{1} Q_n(u)\ du
		= 2M C_n \int_{\delta}^{1} {(1 - u^2)}^n\ du
		\le 2M C_n {(1 - \delta^2)}^n (1 - \delta)
		\le 2M C_n {(1 - \delta^2)}^n
	\]
	נרצה להעריך את $C_n$,
	\[
		C_n \int_{-1}^{1} {(1 - u^2)}^n\ du = 1
	\]
	אז,
	\[
		\int_{-1}^{1} {(1 - u^2)}^n\ du
		\ge \int_{-\frac{1}{\sqrt{n}}}^{\frac{1}{\sqrt{n}}} {(1 - u^2)}^n\ du
		= 2 \int_0^{\frac{1}{\sqrt{n}}} {(1 - u^2)}^n\ du
		\ge 2 \int_0^{\frac{1}{\sqrt{n}}} 1 - u^2\ du
		= 2 {\left[ \frac{1}{\sqrt{n}} - n \frac{u^3}{3}\right]}_0^{\frac{1}{\sqrt{n}}}
	\]
	ולכן נסיק ש־$C_n \le \sqrt{n}$.
	בהתאם נקבל חסם ל־$I_3$ ומטעמי סימטריה גם ל־$I_1$, ונוכל להסיק שלכל $\epsilon > 0$ קיים $\delta > 0$ כך שמתקיים,
	\[
		|F(x) - P_n(x)|
		\le \epsilon + 4M \sqrt{n} {(1 - \delta^2)}^n
		\xrightarrow{n \to \infty} 0
	\]
	כלומר קיים $M_0 > 0$ כל ש־$|F(x) - P_n(x)| \le 2 \epsilon$ לכל $x \in \RR, n > M_0$, ולכן $P_n \overset{\RR}{\rightrightarrows} f$ ובפרט $P_n \overset{[0, 1]}{\rightrightarrows} f$ כפי שרצינו.
\end{proof}

\section{שיעור 5 --- 7.5.2025}
\subsection{קירובים במרחבים מטריים}
ראינו את משפט ויירשטראס לקירוב פונקציות ב־$C([a, b])$.
נניח עתה ש־$(X, \rho)$ מרחב מטרי כלשהו, ונניח ש־$K \subseteq X$, \\
נבחן את $C(K) = \{ f : K \to \RR \mid f \text{ is continuous} \}$. \\
נגדיר את הנורמה $\lVert f \rVert_\infty = \sup_{x \in K} |f(x)|$, ואנו יודעים כי $(C(K), \lVert \cdot \rVert_\infty)$ הוא מרחב נורמי.
המטרה שלנו היא למצוא גרסה כללית יותר של משפט ויירשטראס, הוא משפט סטון־ויירשטראס, כך שתהי $A \subseteq C(K)$ הצפופה ב־$C(K)$.
לשם כך ננסה להכליל את הקונספט של פולינומים.
\begin{definition}[אלגברה]
	נניח ש־$A \subseteq C(K)$ עבור $K \subseteq X$ במרחב המטרי $(X, \rho)$.
	אם התנאים הבאים מתקיימים,
	\begin{enumerate}
		\item אם $f, g \in A$ אז $f + g \in A$
		\item אם $f, g \in A$ אז $f g \in A$
		\item אם $f \in A$ ו־$\alpha \in \RR$ אז $\alpha f \in A$
	\end{enumerate}
	אז נאמר ש־$A$ היא אלגברה.
\end{definition}
\begin{definition}[הפרדת נקודות]
	נניח ש־$A \subseteq C(K)$ אלגברה,
	אם עבור כל $x, y \in K$ כך ש־$x \ne y$ קיימת פונקציה $f \in A$ כך ש־$f(x) \ne f(y)$ אז נאמר ש־$A$ מפרידה נקודות ב־$K$.
\end{definition}
\begin{definition}[אלגברה שאינה מתאפסת באף נקודה]
	נניח ש־$A \subseteq C(K)$,
	אם עבור כל $x \in K$ קיימת פונקציה $f \in A$ כך ש־$f(x) \ne 0$, אז נאמר ש־$A$ אינה מתאפסת באף נקודה.
\end{definition}
עתה נראה מספר דוגמות להגדרות אלה.
\begin{example}
	$A = C(K)$ עבור $K \subseteq \RR$, נבחין כי זוהי אכן אלגברה, כנביעה מהסגירות של מרחב הפונקציות לאלגברה.
	\begin{enumerate}
		\item $A$ מפרידה נקודות,
			זאת שכל לכל $x$ נוכל לבחור את $f(x) = x$
		\item $A$ אינה מתאפסת באף נקודה,
			ההוכחה לזה היא בחירת $f(x) = c$ עבור $c \ne 0$ כלשהו.
	\end{enumerate}
\end{example}
\begin{example}
	הפעם נגדיר את $A = P$ מרחב הפולינומים,
	הפעם גם $A$ מפרידה בין נקודות ואינה מתאפסת באף נקודה.
\end{example}
נעבור לדוגמה נגדית.
\begin{example}
	נגדיר $A \subseteq C[-1, 1]$ כך שמתקיים,
	\[
		A_{\text{even}}
		= \{ f \in C[-1, 1] \mid f \text{ is continuous}, \forall x \in [-1, 1], f(x) = f(-x) \}
	\]
	קבוצת הפונקציות הזוגיות.
	זוהי בבירור אלגברה, שכן מכפלת פונקציות זוגיות היא זוגית וכך גם חיבורן.
	אבל $A_{\text{even}}$ לא מפרידה בין נקודות.
\end{example}
\begin{definition}[קבוצה קומפקטית]
	נניח ש־$(X, \rho)$ מרחב מטרי, ותהי $K \subseteq X$.
	נאמר ש־$K$ קומפקטית אם לכל כיסוי פתוח של $K$ יש תת־כיסוי סופי. \\
	כלומר אם $K \subseteq \bigcup_{\alpha \in I} U_{\alpha}$ עבור קבוצת אינדקסים כלשהי $I$ של קבוצות פתוחות $U_{\alpha}$,
	במקרה זה קיימים $\alpha_1, \ldots, \alpha_N \in I$ כך שמתקיים $K \subseteq \bigcup_{i = 1}^N U_{\alpha_i}$.
\end{definition}
משפט חשוב שמגיע אלינו מטופולוגיה ולא נוכיח במסגרת קורס זה הוא המשפט הבא.
\begin{theorem}[הגדרות שקולות של קומפקטיות]
	יהי $(X, \rho)$ מרחב מטרי ויהי $K \subseteq X$, אז התנאים הבאים שקולים,
	\begin{enumerate}
		\item $K$ קומפקטית
		\item $K$ קומפקטי סדרתית,
			כלומר כל סדרה ב־$K$ מכילה תת־סדרה מתכנסת לנקודה בקבוצה $K$
		\item $K$ שלמה, דהינו חסומה לחלוטין וחסומה
	\end{enumerate}
\end{theorem}
\begin{theorem}[סטון־ויירשטראס]\label{stone_weistrass_theorem}
	נניח ש־$(X, \rho)$ מרחב מטרי,
	$K \subseteq X$ קבוצה קומפקטית המשרה $C(K) = \{ f : K \to \RR \mid f \text{ is continuous} \}$. \\
	נגדיר גם $\lVert f \rVert_\infty = \sup_{x \in K} |f(x)|$.
	במרחב הנורמי $(C(K), \lVert \cdot \rVert_\infty)$. \\
	נניח גם ש־$A \subseteq C(K)$ אלגברה מפרידה בין נקודות ושאינה מתאפסת באף נקודה,
	אז $\overline{A} = C(K)$.
\end{theorem}
לפני שניגש להוכחת המשפט, נגדיר ונוכיח מספר למות.
\begin{lemma}
	נניח ש־$A \subseteq C(K)$ ונניח ש־$A$ אלגברה מפרידה בין נקודות שאינה מתאפסת באף נקודה. \\
	נניח ש־$x, y \in K$ כך ש־$x \ne y$, וכן ש־$c_1, c_2 \in \RR$. \\
	אז קיימת $f \in A$ כך ש־$f(x) = c_1, f(y) = c_2$.
\end{lemma}
\begin{proof}
	קיימות $g, h_1, h_2 \in A$ כך שמתקיים,
	\[
		g(x) \ne g(y),
		\quad h_1(x) \ne 0,
		\quad h_2(y) \ne 0
	\]
	נגדיר את הפונקציות $u(t) = h_2(t) (g(t) - g(y)) \in A$ וכן $v(t) = h(t) (g(t) - g(y)) \in A$, כאשר השייכות ל־$A$ נובעת מהיותה אלגברה.
	מתקיים,
	\[
		u(x) = 0,
		\quad u(y) \ne 0,
		\quad v(x) \ne 0,
		\quad v(y) = 0
	\]
	נגדיר עתה,
	\[
		f(t) = c_1 \frac{v(t)}{v(x)} + c_2 \frac{u(t)}{u(y)}
	\]
	אז מתקיים $f(x) = c_1, f(y) = c_2$.
\end{proof}
נסמן למה זו ב־$(*)$.
\begin{proof}[הוכחת משפט\ \ref{stone_weistrass_theorem}]
	בשלב הראשון יהי $\epsilon > 0$, תהי $f \in C(K)$ ונניח ש־$x \in K$.
	נרצה לבנות פונקציה $g_x$ כך שמתקיים,
	\begin{itemize}
		\item $g_x \in \overline{A}$
		\item $g_x(x) = f(x)$
		\item $g_x(t) > f(t) - \epsilon$ לכל $t \in K$
	\end{itemize}
	עבור כל $y \in K$ קיימת $(*)$ פונקציה $h_y \in A$ כך ש־$h_y(y) = f(y)$ ו־$h_y(x) = f(x)$. \\
	נגדיר את הקבוצה,
	\[
		J_y
		= \{ t \in K \mid h_y(t) > f(t) - \epsilon \}
	\]
	אנו יודעים כי $h_y(y) = f(y) > f(y) - \epsilon$ ולכן $y \in J_y$.
	הקבוצה $J_y$ היא פתוחה מהגדרתה ב־$K$, נוכל לראות זאת מהטופולוגיה המושרית על $K$ מ־$X$.
	נבחין כי הקבוצות $J_y$ מכסות את $K$, כלומר,
	\[
		K
		= \bigcup_{y \in K} J_y
	\]
	כיסוי פתוח של $K$, אבל מקומפקטיות $K$ והאפיון השקול לקומפקטיות במרחבים מטריים נובע שיש תת־כיסוי סופי ל־$K$.
	נסמן,
	\[
		K = \bigcup_{i = 1}^n J_{y_i}
	\]
	עבור $y_i \in K$ לכל $1 \le i \le n$, כאשר $n \in \NN$. \\
	נגדיר $g_x = \max\{ h_{y_1}, \ldots, h_{y_n} \}$ בנורמה $\lVert \cdot \rVert_\infty$.
	נובע ש־$g_x \in \overline{A}$ מ־$(\#)$.
	נבחין גם כי $h_{y_1}(x) = \cdots = h_{y_n}(x)$ ולכן נוכל להסיק שבפרט $g_x(x) = f(x)$.
	לכל $t \in K$ אנו יודעים כי $t \in J_{y_i}$ עבור איזשהו $i$, ולכן,
	\[
		g_x(t)
		\ge h_{y_i}(t)
		> f(t) - \epsilon
	\]
	כאשר קיים $i$ כזה מהעובדה שיש כיסוי סופי.

	בשלב השני נרצה למצוא $\varphi \in \overline{A}$ כך שיתקיים,
	\[
		\lVert \varphi - f \rVert_\infty < \epsilon
	\]
	נגדיר $\hat{J}_x = \{ t \in K \mid g_x(t) < f(t) + \epsilon \}$.
	אבל,
	\[
		g_x(x)
		= f(x)
		< f(x) + \epsilon
	\]
	ולכן $x \in \hat{J}_x$.
	נוכל שוב להגדיר $K = \bigcup_{x \in K} \hat{J}_x$ ושוב קיימים $x_1, \ldots, x_n$ כך שמתקיים,
	\[
		J
		= \bigcup_{i = 1}^n \hat{J}_{x_i}
	\]
	ונגדיר $\varphi = \min\{ g_{x_1}, \ldots, g_{x_n} \}$.
	נעיר כי עלינו להראות ש־$\varphi \in \overline{A}$.
	לכל $t \in K$ קיים $i$ כך ש־$t \in \hat{J}_{x_i}$, ונשים לב שמתקיים,
	\[
		g_{x_i}(t)
		< f(t) + \epsilon
	\]
	ולכן בפרט $\varphi(t) \le g_{x_i}(t) < f(t) + \epsilon$.
	וכן, $\varphi(t) = g_{x_i}(t) > f(t) - \epsilon$.
	נסיק שמתקיים,
	\[
		f(t) - \epsilon
		< \varphi(t)
		< f(t) + \epsilon
	\]
	לכל $t \in K$.
	נובע ש־$|\varphi(t) - f(t)| \le 2 \epsilon$ לכל $t \in K$, ולכן גם $\sup |\varphi(t) - f(t)| \le 2 \epsilon$ לכל $\epsilon > 0$.
\end{proof}

\listoftheorems[title=הגדרות ומשפטים,ignoreall,show={theorem,definition},swapnumber,onlynamed={proposition}]

\end{document}
